\sloppy

% Настройки стиля ГОСТ 7-32
% Для начала определяем, хотим мы или нет, чтобы рисунки и таблицы нумеровались в пределах раздела, или нам нужна сквозная нумерация.
% \EqInChapter % формулы будут нумероваться в пределах раздела
% \TableInChapter % таблицы будут нумероваться в пределах раздела
% \PicInChapter % рисунки будут нумероваться в пределах раздела

% Добавляем гипертекстовое оглавление в PDF
\usepackage[
bookmarks=true, colorlinks=true, unicode=true,
urlcolor=black,linkcolor=black, anchorcolor=black,
citecolor=black, menucolor=black, filecolor=black,
]{hyperref}

\urlstyle{same}
\makeatletter
\def\UrlBreaks{\do\/\do\.\do\-\do\_\do\?\do\&\do\=}
\def\UrlBigBreaks{\do\/}
\makeatother
\newcommand{\BibUrl}[1]{\newline URL: \mbox{\url{#1}}}
\usepackage{csquotes}
\usepackage[
  backend=biber,
  bibstyle=gost-numeric,
  citestyle=gost-numeric,
  sorting=none,
  language=auto,
  autolang=other,
  doi=true,
  url=true,
  eprint=false
]{biblatex}
\DefineBibliographyStrings{russian}{urlseen={дата обращения}}
\DeclareFieldFormat{urldate}{\mkbibparens{\bibstring{urlseen}\addcolon\space#1}}
\addbibresource{rpz.bib}
\renewcommand*{\multicitedelim}{\addcomma\space}
\renewcommand*{\compcitedelim}{\textendash}

\usepackage{pdfpages}

% Изменение начертания шрифта --- после чего выглядит таймсоподобно.
% apt-get install scalable-cyrfonts-tex

% \IfFileExists{cyrtimes.sty}
%     {
        % \usepackage{cyrtimespatched}
%     }
%     {
%         % А если Times нету, то будет CM...
%     }

\usepackage{graphicx}   % Пакет для включения рисунков
\graphicspath{{images/}{image/}}
\usepackage{pdflscape}
\usepackage{eso-pic}

% С такими оно полями оно работает по-умолчанию:
% \RequirePackage[left=20mm,right=10mm,top=20mm,bottom=20mm,headsep=0pt]{geometry}
% Если вас тошнит от поля в 10мм --- увеличивайте до 20-ти, ну и про переплёт не забывайте:

% Пакет Tikz
\usepackage{tikz}
\usetikzlibrary{arrows,positioning,shadows}

% Произвольная нумерация списков.
\usepackage{enumerate}

% ячейки в несколько строчек
\usepackage{multirow}

% itemize внутри tabular
\usepackage{paralist,array}

% Таблицы
\usepackage{array} % расширенные возможности для работы с таблицами
\usepackage{tabularx} % автоматический подбор ширины столбцов
\usepackage{longtable} % многостраничные таблицы
\usepackage{ltablex} % гибрид longtable + tabularx
\usepackage{ragged2e}
\usepackage{makecell}
\usepackage{booktabs}
\usepackage{dcolumn} % выравнивание чисел по разделителю
\keepXColumns
\newcolumntype{Y}{>{\RaggedRight\hyphenpenalty=10000\exhyphenpenalty=10000\arraybackslash}X}
\newcolumntype{P}[1]{>{\RaggedRight\hyphenpenalty=10000\exhyphenpenalty=10000\arraybackslash}p{#1}}
\newcommand{\tableheadcell}[1]{\Centering\hyphenpenalty=10000\exhyphenpenalty=10000\arraybackslash #1}
\renewcommand{\arraystretch}{1}
\setlength{\LTpre}{0pt}
\setlength{\LTpost}{0pt}

\usepackage{setspace} % полуторный интервал
\setstretch{1.5}
\usepackage{indentfirst}
\setlength{\parindent}{1.25cm}

\newcounter{compactlistitem}
\newenvironment{compactlist}{%
    \begin{list}{\arabic{compactlistitem})}{%
        \usecounter{compactlistitem}%
        \setlength{\leftmargin}{0pt}%
        \setlength{\labelwidth}{0pt}%
        \setlength{\labelsep}{0pt}%
        \setlength{\itemindent}{1.25cm}%
        \def\makelabel##1{##1\hspace{0.2cm}}%
        \setlength{\topsep}{0pt}%
        \setlength{\partopsep}{0pt}%
        \setlength{\itemsep}{0pt}%
        \setlength{\parsep}{0pt}%
    }%
}{%
    \end{list}%
}

\makeatletter
\renewcommand*\l@chapter{\@dottedtocline{0}{0em}{1.8em}}
\renewcommand*\l@section{\@dottedtocline{1}{0em}{2.4em}}
\renewcommand*\l@subsection{\@dottedtocline{2}{0em}{3.4em}}
\renewcommand*\l@subsubsection{\@dottedtocline{3}{0em}{4.4em}}
\makeatother
\setcounter{tocdepth}{2}
\setcounter{secnumdepth}{2}


\usepackage{caption}

\DeclareCaptionJustification{centerednohyphen}{\Centering\hyphenpenalty=10000\exhyphenpenalty=10000}
\DeclareCaptionJustification{leftnohyphen}{\RaggedRight\hyphenpenalty=10000\exhyphenpenalty=10000}
\DeclareCaptionLabelSeparator{shortdash}{ \textendash\ }

\usepackage{caption}

\usepackage{float}
\setlength{\floatsep}{0pt}
\setlength{\textfloatsep}{0pt}
\setlength{\intextsep}{7pt}

\usepackage{algorithm}
\usepackage{algpseudocode}
\floatname{algorithm}{Алгоритм}
\algrenewcommand\algorithmicrequire{Входные данные:}
\algrenewcommand\algorithmicensure{Результат:}
\algrenewcommand\algorithmicforall{для каждого}
\algrenewcommand\algorithmicfor{для}
\algrenewcommand\algorithmicdo{выполнить}
\algrenewcommand\algorithmicend{конец}
\algrenewcommand\algorithmicreturn{вернуть}

\captionsetup{aboveskip=7pt,belowskip=7pt,justification=centerednohyphen,singlelinecheck=false,labelsep=shortdash}
\captionsetup[table]{justification=leftnohyphen,singlelinecheck=false,position=top,font={stretch=1},aboveskip=7pt,belowskip=7pt}
\captionsetup[figure]{justification=centerednohyphen,singlelinecheck=false}
\captionsetup[listing]{justification=leftnohyphen,singlelinecheck=false,position=top,font={stretch=1},aboveskip=0pt,belowskip=7pt}
\captionsetup[listing]{position=top}

\usepackage{chngcntr}

\usepackage{mdframed}
\usepackage{xcolor}
\newcommand{\framedincludegraphics}[2][]{%
  {\setlength{\fboxsep}{0pt}\setlength{\fboxrule}{0.4pt}\fbox{\includegraphics[#1]{#2}}}%
}

\mdfsetup{
linecolor=black,
linewidth=1pt,
roundcorner=5pt,
innerleftmargin=5pt,
innerrightmargin=5pt,
innertopmargin=5pt,
innerbottommargin=5pt,
skipabove=0pt, % Reduces space above the frame
skipbelow=0pt, % Reduces space below the frame
leftmargin=0pt, % Reduces left outer margin
rightmargin=0pt % Reduces right outer margin
}

\usepackage[chapter]{minted}
\usepackage{upquote}

% Configure captions for minted's listing environment
\captionsetup[listing]{
    position=top,
    justification=leftnohyphen,
    singlelinecheck=false,
    font={stretch=1},
    aboveskip=0pt,
    belowskip=7pt
}

% Alternative: Force all float captions to top
\floatstyle{plaintop}
\restylefloat{listing}

\usemintedstyle{bw}
\setminted{
    style=bw,
    linenos=true,           % включить номера строк
    numberblanklines=false, % не нумеровать пустые строки
    numbersep=-16pt,
    xleftmargin=0pt,
    fontsize=\small,
    fontfamily=listingfont,
    baselinestretch=1,
    frame=single,
    framesep=24pt,
    framerule=0.4pt,
    rulecolor=\color{black},
    breakanywhere=true,
    breakautoindent=true,
    breaklines=true,
}

\newenvironment{listingblock}[1][35mm]{%
    \par\addvspace{7pt}%
    \Needspace{#1}%
    \noindent
}{%
    \par
}

\newenvironment{appendixlistingblock}[1][24mm]{%
    \par\addvspace{7pt}%
    \Needspace{#1}%
    \noindent
}{%
    \par
}

\newcommand{\listingpartlabel}[1]{%
    \noindent{\normalfont\fontsize{14pt}{21pt}\selectfont #1}\par\nobreak\addvspace{7pt}%
}

\newcommand{\tablepartlabel}[1]{%
    {\normalfont\fontsize{14pt}{21pt}\selectfont #1}%
}

\newcommand{\tablecontinuedorend}[2]{%
    \multicolumn{#1}{l}{\tablepartlabel{\ifnum\value{page}=#2 окончание\else продолжение\fi\ таблицы \thetable}}\\%
}

\renewcommand{\listingscaption}{Листинг}
\renewcommand{\listoflistingscaption}{Листинги}


\makeatletter
\let\oldtableofcontents\tableofcontents
\renewcommand{\tableofcontents}{%
  \begingroup
  \hyphenpenalty=10000
  \exhyphenpenalty=10000
  \oldtableofcontents
  \endgroup
}
\makeatother

\makeatletter
\def\@hangfrom#1{\setbox\@tempboxa\hbox{{#1}}%
      \hangindent 0pt%\wd\@tempboxa
      \noindent\box\@tempboxa}
\makeatother

\usepackage{etoolbox}
\makeatletter
\AtBeginEnvironment{minted}{%
  \let\textbf\@firstofone
  \let\textit\@firstofone
  \let\textsl\@firstofone
  \let\emph\@firstofone
}
\makeatother
\AfterEndEnvironment{minted}{\par\addvspace{7pt}}
\AtBeginEnvironment{longtable}{\small\setstretch{1}\selectfont}
\AtEndEnvironment{listing}{\vspace{0pt}}
\AtEndEnvironment{table}{\vspace{0pt}}
\AtEndEnvironment{center}{\vspace{0pt}}
\AtEndEnvironment{figure}{\vspace{0pt}}


\usepackage{fontspec}

% Указываем шрифты для основного текста и для кириллицы:
\IfFontExistsTF{Times New Roman}{
  \setmainfont{Times New Roman}
  \setsansfont{Times New Roman}
  \setmonofont{Times New Roman}[Mapping=]
  \newfontfamily\cyrillicfont{Times New Roman}
}{
  \PackageError{preamble}{Times New Roman is required but not installed}
    {Install Times New Roman and rebuild with XeLaTeX. Do not use Liberation Serif as a substitute.}
}
\IfFontExistsTF{Courier New}{
  \newfontfamily\listingfont{Courier New}[
    Scale=MatchLowercase,
    BoldFont={Courier New},
    ItalicFont={Courier New},
    BoldItalicFont={Courier New}
  ]
}{
  \newfontfamily\listingfont{DejaVu Sans Mono}[
    Scale=MatchLowercase,
    BoldFont={DejaVu Sans Mono},
    ItalicFont={DejaVu Sans Mono},
    BoldItalicFont={DejaVu Sans Mono}
  ]
}
\makeatletter
\@namedef{FV@fontfamily@listingfont}{%
  \def\FV@FontScanPrep{}%
  \let\FV@FontFamily\listingfont
}
\makeatother
\renewcommand{\theFancyVerbLine}{%
  \normalfont\listingfont\fontsize{12pt}{12pt}\selectfont\arabic{FancyVerbLine}%
}

\makeatletter
\setlength\GostAfterTitleSkip{14pt}
\setlength\GostBeforTitleSkip{14pt}
\renewcommand\GostTitleStyle{\normalfont\bfseries\raggedright\hyphenpenalty=10000\exhyphenpenalty=10000}
\renewcommand\@eqnnum{\normalfont\normalcolor(\theequation)}
\newcommand*\LandscapePageNumber{%
  \AtPageLowerLeft{%
    \put(\LenToUnit{\dimexpr\paperwidth-10mm\relax},\LenToUnit{0.5\paperheight}){%
      \makebox(0,0){\rotatebox{90}{\normalfont\fontsize{14pt}{14pt}\selectfont\thepage}}%
    }%
  }%
}
\let\VKR@landscape\landscape
\let\endVKR@landscape\endlandscape
\renewenvironment{landscape}{%
  \clearpage
  \pagestyle{empty}%
  \AddToShipoutPictureFG{\LandscapePageNumber}%
  \VKR@landscape
}{%
  \endVKR@landscape
  \ClearShipoutPictureFG
  \pagestyle{plain}%
}
\makeatother

% Убирает точку после номера главы
\makeatletter
\renewcommand{\thesection}{\thechapter.\arabic{section}}
\renewcommand{\thesubsection}{\thesection.\arabic{subsection}}
\renewcommand{\@seccntformat}[1]{%
  {\fontsize{14pt}{21pt}\selectfont\bfseries \csname the#1\endcsname}\hspace{1em}%
}
\makeatother
